\chapter{Method AV: artificial viscosity}
\label{ch:method-av}

\checkedagainst{1.9.0 (\code{11ecf43})}{Each update against
\code{src/modules/av_mod.f95} and the cycle in \code{maincode.f95},
section~7. Conservation is checked exactly (rational arithmetic) and the
energy update's split symbolically in \code{doc/manual/checks/ch4_av.py},
which also measures the shock width quoted here.}

Method AV is \bw's default (deck: \key{method = AV}), and the ME~697
class project's. It is the scheme of von Neumann and Richtmyer: velocity
and position on the faces of the cells, the thermodynamic state at their
centers, a leapfrog in time, and an artificial viscosity $q$ that spreads
each shock over a few cells. This chapter derives each update from the
Lagrangian equations of \cref{ch:continuum} and the shock structure of
\cref{ch:shocks}, and says what the code does where it goes beyond the
course's sheet (\code{doc/Hydrocode_Eqns_v4}). Its material models (the
pressure, tension and spall, the sound speed) are shared with method RS
and are in \cref{ch:materials}.

\section{The staggered mesh and the leapfrog}
\label{sec:av-mesh}

Cell $k$ holds the state that belongs to a piece of material: its density
$\rho_k$, internal energy $e_k$, pressure $P_k$, viscosity $q_k$ and sound
speed $c_k$, and its fixed mass $M_k$ (\cref{eq:cm-mass-cell}). Its faces,
$k - \tfrac{1}{2}$ and $k + \tfrac{1}{2}$, hold positions $x$ and
velocities $u$ (\cref{fig:av-grid}). A face moves as the material on its
two sides pushes it, so it carries the mass of half of each neighbor, the
\emph{node mass}
\begin{equation}
m_{k+1/2} = \tfrac{1}{2}\bigl(M_k + M_{k+1}\bigr).
\label{eq:av-node-mass}
\end{equation}
This is the staggering: each update reads quantities that sit where it
needs them, so every difference in it is centered in space.

\begin{figure}[ht]
\centering
\includegraphics[width=0.7\textwidth]{figures/hc_grid.pdf}
\caption{The staggered mesh. Velocity and position live on the faces; the
thermodynamic state at the cell centers. A face's momentum update divides
by its node mass, half of each neighbor cell.}
\label{fig:av-grid}
\end{figure}

The time levels are staggered too (\cref{fig:av-leapfrog}). The state of a
cell is known at the whole steps $n$, $n + 1$; the face velocities at the
half steps $n - \tfrac{1}{2}$, $n + \tfrac{1}{2}$. Each update is centered
on the level of what drives it: the velocity jumps from $n - \tfrac{1}{2}$
to $n + \tfrac{1}{2}$ under the stresses at $n$, and the positions and the
state go from $n$ to $n + 1$ with the velocity at $n + \tfrac{1}{2}$.

\begin{figure}[ht]
\centering
\includegraphics[width=0.6\textwidth]{figures/hc_leapfrog.pdf}
\caption{The leapfrog. Velocity is advanced on half steps by the state at
the whole step between them; everything else is advanced on whole steps
by the new velocity.}
\label{fig:av-leapfrog}
\end{figure}

\paragraph{The steps.} $\Delta t^n$ is the step from $n$ to $n + 1$. The
velocity's step, from $n - \tfrac{1}{2}$ to $n + \tfrac{1}{2}$, is the
mean of the previous step and the current one,
\begin{equation}
\Delta t^{n-1/2} = \tfrac{1}{2}\bigl(\Delta t^{n-1} + \Delta t^n\bigr),
\label{eq:av-dth}
\end{equation}
which is also the time between the two half steps, since each half step
lies midway through its whole step. This keeps the velocity update centered when the step changes from cycle
to cycle (\code{dth} in the code). At the start, the velocity at
$t = 0$ stands in for the one at $-\tfrac{1}{2}$, and $\Delta t^{-1} =
\Delta t^0$: the first step is first order, every later one centered.

\section{Momentum and the moving mesh}
\label{sec:av-momentum}

Apply the momentum balance of a slab, \cref{eq:cm-cell}, to the slab
between the centers of cells $k$ and $k + 1$: its mass is the node mass,
its velocity the face's, and the stresses on it are those at the two
centers, $\sigma = -(P + q)$. Over the velocity's step,
\begin{equation}
u^{n+1/2}_{k+1/2} = u^{n-1/2}_{k+1/2}
  + \frac{\Delta t^{n-1/2}}{m_{k+1/2}}
    \Bigl[\bigl(P^n_k - P^n_{k+1}\bigr) + \bigl(q^n_k - q^n_{k+1}\bigr)\Bigr].
\label{eq:av-momentum}
\end{equation}
\trace{modules/av_mod.f95}{face_step}{compare_g1.py, study_av.py}
The face then moves with the material, and the cell's density follows from
its fixed mass:
\begin{equation}
x^{n+1}_{k+1/2} = x^n_{k+1/2} + u^{n+1/2}_{k+1/2}\,\Delta t^n,
\qquad
\rho^{n+1}_k = \frac{M_k}{x^{n+1}_{k+1/2} - x^{n+1}_{k-1/2}} .
\label{eq:av-position}
\end{equation}
\trace{modules/av_mod.f95}{face_step, cell_step}{compare_g1.py}
These are the Lagrangian equations of momentum and mass,
\cref{eq:lag-mom,eq:lag-mass}, in difference form. In the mass coordinate
the centers of cells $k$ and $k + 1$ are $m_{k+1/2}$ apart, so
\cref{eq:av-momentum} is $\partial u/\partial t = \partial\sigma_{xx}/\partial
h$ with $\sigma_{xx} = -(P + q)$, differenced across a face. A cell's faces
are $M_k$ apart, and its specific volume is its width over its mass, so
\cref{eq:av-position} gives $(V^{n+1}_k - V^n_k)/\Delta t^n = (u^{n+1/2}_{k+1/2}
- u^{n+1/2}_{k-1/2})/M_k$, which is $\partial V/\partial t = \partial
u/\partial h$ differenced across a cell. Each cell's $P + q$ pushes its two
faces equally in opposite directions, so the sum of $m_{k+1/2}u_{k+1/2}$
over all faces, the stack's momentum, changes only through the two ends,
where the void pushes with zero stress.

\section{Energy}
\label{sec:av-energy}

The energy update is where the scheme is easiest to get wrong. \bw's is
built to conserve \emph{total} energy exactly, from the cell form of
\cref{eq:cm-cell}: a cell's total energy changes only by the work done on
its two faces,
\begin{multline}
e^{n+1}_k = e^n_k + \frac{\Delta t^n}{M_k}\Bigl[\,\bar\sigma_{k-1/2}\,
  u^{n+1/2}_{k-1/2} - \bar\sigma_{k+1/2}\,u^{n+1/2}_{k+1/2}\Bigr] \\
  + \tfrac{1}{8}\Bigl[\bigl(u^{n-1/2}_{k-1/2} + u^{n-1/2}_{k+1/2}\bigr)^2
  - \bigl(u^{n+1/2}_{k-1/2} + u^{n+1/2}_{k+1/2}\bigr)^2\Bigr].
\label{eq:av-energy}
\end{multline}
\trace{modules/av_mod.f95}{cell_step}{compare_g1.py, study_av.py}
The first bracket is the work of the face stresses over the step. The
second takes away the change in the cell's kinetic energy,
$\tfrac{1}{2}\bar u^2$, with $\bar u$ the mean of its two face velocities.
The stress on a face, positive in compression, is the pressure interpolated
linearly in mass between the two cell centers, plus their average
viscosity:
\begin{equation}
\bar\sigma_{k+1/2} = \frac{M_{k+1}P^n_k + M_k P^n_{k+1}}{M_k + M_{k+1}}
  + \tfrac{1}{2}\bigl(q^n_k + q^n_{k+1}\bigr).
\label{eq:av-face-stress}
\end{equation}
Both cells on a face use the same $\bar\sigma$, so a face's work leaves
one and enters the other: summed over the stack, the work telescopes to
the two ends, and
\begin{equation}
\sum_k M_k\Bigl(e^{n+1}_k + \tfrac{1}{2}\bigl(\bar u^{n+1/2}_k\bigr)^2\Bigr)
= \sum_k M_k\Bigl(e^n_k + \tfrac{1}{2}\bigl(\bar u^{n-1/2}_k\bigr)^2\Bigr)
\label{eq:av-conserve}
\end{equation}
exactly, whatever the momentum update did. This total is the one the
heartbeat reports. The check script confirms it in exact rational
arithmetic.

\paragraph{The end faces.} Beyond each end of the stack sits a ghost cell
of zero mass and zero pressure (\cref{fig:av-void}). Its zero mass makes
the mass-weighted pressure on an end face zero by itself; its viscosity is
not averaged in. An end face carries no $q$ in \cref{eq:av-energy}: the
void does no work.\footnote{The course's sheet averages $q$ with the
ghost's zero, which puts half the end cell's $q$ on the end face. Total
energy was then lost each time a shock reached a free surface; \bw\ has
given the end face no $q$ since 1.3.2 (\code{doc/README.md}, ``Where the
code has moved past a sheet'').} The momentum update needs no special case:
the end face's node mass is half a cell, and the void pushes with zero
stress.

\begin{figure}[ht]
\centering
\includegraphics[width=0.85\textwidth]{figures/hc_void.pdf}
\caption{The ends of the stack: a ghost cell of zero mass, pressure and
viscosity beyond each end, so the end faces are free surfaces.}
\label{fig:av-void}
\end{figure}

\paragraph{What the update does to one cell.} Write $a$ and $b$ for a
cell's left and right face velocities, $\bar\sigma_l$ and $\bar\sigma_r$
for its face stresses, and $\Delta V = (b^{n+1/2} - a^{n+1/2})\,\Delta
t^n/M_k$ for its change of volume. Then \cref{eq:av-energy} is exactly
\begin{equation}
e^{n+1}_k - e^n_k = -\tfrac{1}{2}\bigl(\bar\sigma_l + \bar\sigma_r\bigr)\,\Delta V
  + R_k,
\qquad
R_k = -\bar u^{n+1/2}\Bigl[\frac{\Delta t^n}{M_k}\bigl(\bar\sigma_r -
  \bar\sigma_l\bigr) + \Delta\bar u\Bigr] + \tfrac{1}{2}\Delta\bar u^2,
\label{eq:av-energy-split}
\end{equation}
with $\Delta\bar u = \bar u^{n+1/2} - \bar u^{n-1/2}$. The first term is the
first law of \cref{eq:lag-energy}, $\dd e = -(P + q)\,\dd V$. The
remainder $R_k$ is the price of exact conservation. The mesh moves a cell's
faces with their own node momenta, not with the cell's mean, so the bracket
is not exactly zero; if it were, $R_k$ would be $\tfrac{1}{2}\Delta\bar u^2$,
second order in the step. In smooth flow $R_k$ is small. It is not small in
one place: the cell at a free surface, when a release blows it off, where
$\bar u$ grows fast and its own $-P\,\dd V$ is near zero. There the growth of
$\tfrac{1}{2}\bar u^2$ is taken from $e$, and $e$ can fall well below what
the cell's compression history would give it. A cell's temperature is
computed from its internal energy (\cref{ch:materials}), so this is where
the error shows. In the one measured case, a barely burnt end cell of
Material X blown off a free surface, the computed temperature reaches about
$-1000$~K on every mesh, in that one cell of vanishing mass, while total
energy stays exact
(\code{doc/design/av_study.md}, set C). A ``compatible'' energy update, in
which each cell is charged exactly $-(P + q)\,\dd V$, would remove it, at the
cost of the course's scheme (Caramana et~al.\ \cite{caramana1998}).

\paragraph{Time levels.} The stresses in \cref{eq:av-energy} are those at
$n$, the start of the step, with the velocities at $n + \tfrac{1}{2}$.
Total energy is conserved exactly, but the pressure work is taken at the
step's start: first order in time. The order of accuracy of the whole
scheme on smooth flow has not yet been measured in a test (\code{TODO.md},
section~8).

\paragraph{Pressure.} With $\rho^{n+1}_k$ and $e^{n+1}_k$ known, the
equation of state gives $P^{n+1}_k$, then the spall rule decides whether
the cell holds that pressure or breaks, and the sound speed $c^{n+1}_k$ is
taken at the pressure kept (\cref{ch:materials}). A cell that has started to
burn takes them from the closure of \cref{ch:burning-cell} instead.

\section{The artificial viscosity}
\label{sec:av-q}

\Cref{sec:inside-front} showed that any viscous stress that vanishes
outside a shock, and is positive inside it while the material compresses,
carries the shock between the right end states, with the right entropy
jump, in equations that conserve mass, momentum and energy. Method AV's
is
\begin{equation}
q^{n+1}_k =
\begin{cases}
B_L\,\rho^{n+1}_k c^{n+1}_k\,|\Delta u_k| + B_Q\,\rho^{n+1}_k\,(\Delta u_k)^2
  & \text{if } \Delta u_k > 0 \text{ and } \rho^{n+1}_k > \rho^n_k, \\
0 & \text{otherwise,}
\end{cases}
\label{eq:av-q}
\end{equation}
\trace{modules/av_mod.f95}{cell_state}{compare_g1.py, study_av.py}
with $\Delta u_k = u^{n+1/2}_{k-1/2} - u^{n+1/2}_{k+1/2}$ the cell's
velocity jump, positive while it is squeezed, and $B_L$ and $B_Q$ deck keys
(\key{BL}, \key{BQ}; 0.1 and 2.0 by default).

\begin{itemize}
\item \textbf{The quadratic term} is von Neumann and Richtmyer's
  \cite{vonneumann1950}. Within a
  smeared shock $\Delta u \approx \Delta x\,|\partial u/\partial x|$, so it
  acts as a viscosity $B_Q\,\rho\,\Delta x^2|\partial u/\partial x|$ that
  grows with the velocity gradient and scales with the cell: it spreads
  every shock over a few cells, whatever its strength, and vanishes as the
  mesh refines. Measured in \bw\ (aluminum, $u_p$ from 0.5 to 2~km/s, 25 and
  100 cells per mm), a shock's velocity rises from 10\% to 90\% of its jump
  across 2.1 to 3.2 cells with the default $B_Q = 2$; 1.6 to 2.8 with $B_Q =
  1$; 2.8 to 3.7 with $B_Q = 4$ (\code{checks/ch4_av.py}).
\item \textbf{The linear term} damps the ringing a quadratic viscosity
  alone leaves behind a shock, where $\Delta u$ is small.
\item \textbf{The switch} is the second law of \cref{sec:entropy}: $q$ is
  nonzero only while the cell is being compressed, so $-q\,\dd V \ge 0$
  and it only ever makes entropy. In one dimension the two tests in
  \cref{eq:av-q} agree (a cell's width changes by $-\Delta u_k\,\Delta t^n$);
  the code checks both, as the course's sheet does.
\end{itemize}

The $q$ of \cref{eq:av-q} is at $n + 1$, ready for the next cycle's
momentum and energy updates, in which it acts at $n$.

\section{The time step}
\label{sec:av-dt}

An explicit scheme is stable only if no wave crosses more than one cell in
a step. Each cell's limit is the time for a weak mechanical disturbance (a
sound wave), carried with the flow, to cross it,
\begin{equation}
\Delta t_k = \frac{x^{n+1}_{k+1/2} - x^{n+1}_{k-1/2}}{|\bar u^{n+1/2}_k| + c^{n+1}_k},
\qquad
\Delta t^{n+1} = \mathrm{CFL}\,\min_k \Delta t_k ,
\label{eq:av-cfl}
\end{equation}
\trace{modules/av_mod.f95}{cell_dt}{compare_g1.py}
with CFL a deck key (\key{CFL}, 0.6 by default). In a Lagrangian mesh the
flow does not cross cells, so $|\bar u|$ only shortens the step. The
limit does not include $q$; the default safety factor has been enough for
every deck tested.
A void cell, stretched past the point where its sound speed vanishes
(\cref{ch:materials}), limits the step by its own speed alone,
$\Delta x/|\bar u|$.

\section{One cycle}
\label{sec:av-cycle}

\bw\ runs a cycle of method AV as three loops over the mesh
(\code{maincode.f95}, section~7). Each loop finishes before the next one
begins. A deck with a burn or a temperature adds steps between and after
the loops, in the order listed:

\begin{enumerate}
\item \textbf{Faces:} $u^{n+1/2}$ and $x^{n+1}$, \cref{eq:av-momentum,eq:av-position}.
\item \textbf{Cells:} $\rho^{n+1}$ and $e^{n+1}$, \cref{eq:av-position,eq:av-energy}.
\item \textbf{The burn,} for cells that can react: each model's step,
  from the pressure still at $n$ (\cref{sec:one-cycle}).
\item \textbf{Cells:} $P^{n+1}$ with the spall rule, $c^{n+1}$, $q^{n+1}$
  and $\Delta t_k$, \cref{eq:av-q,eq:av-cfl}.
\item \textbf{Then,} after the loops: the clock advances, the next step
  is set by \cref{eq:av-cfl}, each cell that carries a temperature updates
  it (\cref{ch:materials}), and the burn's lead-shock detector looks at the
  new state (\cref{sec:lead-shock}).
\end{enumerate}

Each loop reads only what the loop before it wrote, and each place's update
reads only its neighbors, so the loops run in parallel with OpenMP and give
the serial answer bit for bit, on any number of threads
(\cref{ch:verification}).

\section{What method AV does, measured}
\label{sec:av-measured}

Method AV converges, but on a coarse mesh it has two known errors, both
measured against exact answers and method RS in
\code{doc/design/av_study.md} (\code{tests/verification/study_av.py}):

\begin{itemize}
\item \textbf{At a free surface,} the velocity overshoots by about 6\% when
  a shock arrives, for about 2.5 cells' transit of the shock (16~ns at 25
  cells per mm, 1~ns at 400), then settles on the right value. The end face
  has half a cell's mass and takes the end cell's full $P + q$, so it rings
  once. Read the plateau, not the peak. (Method RS at order 2: 0.4\%.)
\item \textbf{At a detonation,} AV burns inside its smeared shock, and on a
  coarse mesh the front runs as a weak detonation whose speed is too high:
  $D$ is 3\%
  above $D_{\mathrm{CJ}}$ with 3.5 cells across the reaction zone (25 per mm
  in Material X), 0.4\% with 7, 0.2\% with 14, and within 0.03\% from 28.
  Method RS at order 2 is within 0.02\% at every mesh. The deck reference
  gives the resolution to use (\code{input/README.md}, Material block).
\end{itemize}

Neither is engineered away: each converges as the mesh refines, and each is
documented where a user will look.
